% TOP_PROBLEM: {"id": 30006711, "title": "Woodins Ultimate L Conjecture for Supercompact Cardinals", "category_id": 10, "category": {"id": 10, "name": "set_theory", "display_name": "Set Theory", "description": "Foundations of mathematics, infinite sets, and cardinality.", "slug": "set-theory", "order_index": 10, "created_at": "2026-07-31T15:26:25.671Z"}, "status": "open"}
% FIRST_POSTED: 2026-09-27
% !TeX program = lualatex
% Compile twice: lualatex 30006711.tex
% Standalone research notebook; original section preserved.
% Research additions: 27 September 2026. No external TeX input or BibTeX file.
\documentclass[11pt,a4paper]{article}
\usepackage[margin=24mm,headheight=14pt]{geometry}
\usepackage{fontspec}
\setmainfont{Latin Modern Roman}
\setsansfont{Latin Modern Sans}
\setmonofont{Latin Modern Mono}
\usepackage{amsmath,amssymb,mathtools}
\usepackage{unicode-math}
\setmathfont{Latin Modern Math}
\usepackage{microtype}
\usepackage{enumitem,xurl,needspace}
\usepackage[dvipsnames]{xcolor}
\usepackage{hyperref}
\hypersetup{unicode=true,colorlinks=true,linkcolor=MidnightBlue,urlcolor=MidnightBlue,
 pdftitle={Research notebook - Problem 30006711},pdfauthor={Open Problems Math},bookmarksnumbered=true}
\usepackage{bookmark}
\usepackage{titlesec}
\titleformat{\section}{\Large\bfseries\raggedright}{\thesection}{0.7em}{}
\usepackage{fancyhdr}
\pagestyle{fancy}\fancyhf{}
\fancyhead[L]{\small\sffamily Open Problems Math | Research notebook}
\fancyhead[R]{\small\sffamily Problem 30006711 | 27 September 2026}
\fancyfoot[C]{\thepage}
\setlength{\parindent}{0pt}
\setlength{\parskip}{5pt plus 1pt}
\setlength{\emergencystretch}{3em}
\setcounter{tocdepth}{1}
\setcounter{secnumdepth}{1}
\setlist[itemize]{leftmargin=3em,itemsep=5pt,topsep=4pt,parsep=0pt}
\allowdisplaybreaks
\numberwithin{equation}{section}
\newcommand{\N}{\mathbb{N}}
\newcommand{\Z}{\mathbb{Z}}
\newcommand{\Q}{\mathbb{Q}}
\newcommand{\R}{\mathbb{R}}
\newcommand{\C}{\mathbb{C}}
\newcommand{\F}{\mathbb{F}}
\newcommand{\A}{\mathbb{A}}
\newcommand{\PP}{\mathbb P}
\newcommand{\E}{\mathbb{E}}
\newcommand{\eps}{\varepsilon}
\newcommand{\dd}{\,\mathrm d}
\newcommand{\sref}[2]{\hyperlink{source:#1:#2}{[S#2]}}
\newcommand{\eref}[2]{\hyperlink{extra:#1:#2}{[E#2]}}
\newif\ifshowcatalogue
\showcataloguetrue
\newcommand{\cataloguescope}[1]{\ifshowcatalogue
 \paragraph{Exact scope in the supplied catalogue.}
 \begin{quote}\small #1\end{quote}\fi}
\begin{document}
\setcounter{section}{0}
\section*{\texorpdfstring{Woodins Ultimate L Conjecture for Supercompact Cardinals}{Woodins Ultimate L Conjecture for Supercompact Cardinals}}\label{30006711}
\subsection{Definitions and mathematical statement}
An inner model is a transitive class containing all ordinals and satisfying the specified set theory; $\mathrm{HOD}$ is the class of hereditarily ordinal-definable sets. A cardinal $\delta$ is extendible if sufficiently high rank segments admit elementary embeddings with critical point $\delta$ and image of $\delta$ above the chosen segment. A weak extender model for its supercompactness has, for every $\gamma>\delta$, a normal fine $\delta$-complete ultrafilter $U$ on $\mathcal P_\delta(\gamma)$ with
\[
 \mathcal P_\delta(\gamma)\cap N\in U,\qquad U\cap N\in N.
\]
The selected assertion is the existence of such an $N\subseteq\mathrm{HOD}$ satisfying $V=\mathrm{Ultimate}\text{-}L$. The latter is Woodin's source-defined axiom, not ordinary $V=L$; its full internal-model clauses are incorporated by reference to the cited Conjecture 5.2.
\par\smallskip\textit{Formulation and concept references:} \sref{30006711}{1}.
\cataloguescope{Assuming delta is an extendible cardinal, prove that there is an inner model N contained in HOD such that N is a weak extender model for the supercompactness of delta and N satisfies V=Ultimate-L.}
\subsection{Short English statement}
Does an extendible cardinal guarantee the predicted Ultimate-$L$ inner model inside HOD with the specified supercompactness structure?
% BEGIN RESEARCH_ADDITION ATTEMPT_75
\subsection{Research attempt}

\paragraph{Current frontier and incorporated axiom.}
The source's Conjecture~5.2 was inspected, with its preceding weak-extender
and Ultimate-$L$ conventions \sref{30006711}{1}. In the relevant axiom,
Ultimate-$L$ involves a proper class of Woodin cardinals and reflection of
true $\Sigma_2$ sentences into models $\mathrm{HOD}^{L(A,\R)}$ for suitable
universally Baire $A$; these clauses are interpreted \emph{inside $N$},
not substituted by $N=L$ or by a statement about the ambient continuum.
Aguilera--Bagaria--L\"ucke give a conditional obstruction using an exacting
cardinal above an extendible cardinal \eref{30006711}{1}. Their joint existence
or joint consistency hypothesis is stronger than the catalogue premise.
Neither it nor an unconditional disproof is established in this notebook.

\paragraph{Attack route.}
The constructive route would select coherent normal fine complete measures
and a canonical inner model inside HOD. It must also establish the internal
Ultimate-$L$ axiom, not just supercompactness. A counterexample route can
stop earlier: prevent any weak extender model inside HOD, regardless of
its additional axiom. The latter route has a precise current obstruction,
which is analyzed conditionally below. The former is stress-tested by an
explicit failure of a proposed compactness argument for complete measures.
This separates new notebook reasoning from the deep theorems being used.

\paragraph{Attempt: a sufficient obstruction before Ultimate-$L$.}
For an extendible $\delta$, let $\mathcal H$ denote the HOD Hypothesis in
the source convention: a proper class of regular cardinals are not
$\omega$-strongly measurable in HOD. The relevant established bridge,
attributed to Woodin in \eref{30006711}{1}, Theorems~1.1--1.2, is
\begin{equation}\label{eq75:bridge}
 \exists N\subseteq\mathrm{HOD}\text{ weak extender model for }
 \delta\quad\Longleftrightarrow\quad\mathcal H,
\end{equation}
and the HOD dichotomy makes $\mathcal H$ imply that every singular cardinal
above $\delta$ remains singular in HOD, with correctly computed successor.
Here ``weak extender model'' retains the measure concentration and trace
conditions in the supplied statement; \eqref{eq75:bridge} is not a theorem
about arbitrary inner models merely believing that $\delta$ is supercompact.

Consequently the following is an immediate, but useful, conditional test:
\begin{equation}\label{eq75:obstruction}
 \begin{gathered}
 \delta\text{ extendible},\quad \lambda>\delta\text{ singular in }V,
 \quad \lambda\text{ regular in HOD}\\
 \Longrightarrow\quad
 \text{no such weak extender }N\subseteq\mathrm{HOD}.
 \end{gathered}
\end{equation}
To verify the implication, assume $N$ exists. The forward direction of
\eqref{eq75:bridge} gives $\mathcal H$, and the close-HOD alternative then
makes $\lambda$ singular in HOD, contradicting the last hypothesis.
Thus \eqref{eq75:obstruction} would refute the requested existence conclusion
for that particular $\delta$ before any analysis of Ultimate-$L$ inside $N$.
It does not require HOD itself to satisfy or fail Ultimate-$L$.

The paper defines exacting cardinals by elementary embeddings from suitable
elementary submodels of rank segments. Its Definition~2.4 and Theorem~2.10
supply a source of the singular/regular discrepancy in
\eqref{eq75:obstruction}; Conclusion~6.8 records the resulting conditional
failure. \eref{30006711}{1}
The inference for this catalogue is therefore under the theory
\[
 T=\mathsf{ZFC}+\text{``an exacting cardinal exists above an
 extendible cardinal.''}
\]
No step here derives that extra cardinal from the extendible cardinal
alone. The composition of the source theorems is not claimed as a new
large-cardinal theorem.

\paragraph{Attempt: disentangling model existence from consistency.}
There are three different conclusions one might incorrectly collapse.
If one is actually working in a universe satisfying $T$ and the cited
bridge, \eqref{eq75:obstruction} denies the requested $N$ there. If only
$\operatorname{Con}(T)$ is assumed in an external metatheory, completeness
gives a model of $T$ and hence a conditional model-theoretic obstruction
to proving the universal assertion from the weaker extendible theory.
It does not automatically give a well-founded transitive set model that
can be exhibited as a counterexample within the ambient universe. Finally,
individual consistency assertions for ``extendible'' and ``exacting''
do not imply consistency of their conjunction with the specified ordering.

The elementary logical point can be seen without set theory: a theory can
be consistent with $P$ and separately with $\neg P$ without being
consistent with both. Relative consistency reductions preserve their whole
antecedent; discarding one conjunct destroys the inference. Likewise a
claim that a sufficiently strong choiceless theory yields
$\operatorname{Con}(T)$ still has that stronger theory as a hypothesis.
These qualifications are especially important here because a syntactic
claim that a theory proves a statement and the statement's truth in a
given universe are not interchangeable. The exact conditional outcome is
kept rather than announced as an unconditional solution of the catalogue.

\paragraph{Attempt: a compactness proposal and a concrete failure.}
Consider a different attempt to construct $N$. One might try to satisfy
finitely many measure requirements, then take a limit in a compact space
of ultrafilters and assert that the limiting measures remain complete.
This is false without additional control, even in a much simpler setting.
The set of ultrafilters on $\N$ is a closed subset of
$\{0,1\}^{\mathcal P(\N)}$: the finite Boolean ultrafilter axioms are
closed conditions. Let $U_n$ be the principal ultrafilter at $n$.
Each $U_n$ is closed under arbitrary intersections. Compactness yields a
cluster subnet with indices tending to infinity. Every cofinite subset
belongs eventually to $U_n$, so belongs to the limit $U$.
No finite subset belongs to $U$, and
\[
 \{m:m\geq k\}\in U\quad(k\in\N),
 \qquad \bigcap_{k\geq1}\{m:m\geq k\}=\varnothing\notin U.
\]
Thus a limit of complete principal ultrafilters need not be countably
complete. This argument uses a subnet, not the false assertion that the
whole sequence converges on every subset of $\N$.

The weak-extender requirements involve $\delta$-completeness, normality,
fineness and compatibility with $N$. They are far stronger than the finite
Boolean axioms retained by this compactness argument. A successful
construction must justify preservation of those infinitary requirements,
as well as the fact that $U\cap N$ belongs to $N$. A measure being
ordinal-definable externally would not on its own certify that trace
condition. Selecting a least candidate using HOD's well-order also only
works if an eligible candidate is already an element of the class being
well-ordered; it cannot turn an arbitrary external witness into a HOD
member.

\paragraph{Stress test.}
Neither the ultrafilter example nor the conditional singular-cardinal
obstruction constructs a new inner model or supplies the missing joint
consistency proof. The former refutes only a proposed preservation step;
the latter uses deep source results and an extra hypothesis. Even proving
$N=\mathrm{HOD}$ were a weak extender model would not establish the
separately required internal Ultimate-$L$ clauses. Conversely, an internal
statement such as CH in $N$ need not be CH in $V$, so forcing a different
ambient continuum does not alone contradict this target.

The unverified proof claim in \sref{30006711}{4} is not used as a theorem. No
finite computation can audit the relevant class embeddings or consistency
claims; the companion script explicitly lists this entry as an analytic
and logical argument with no computational certification. Any proposed
stronger conclusion here would have to identify precisely which additional
large-cardinal assumption has been removed and prove that removal.

\paragraph{Outcome.}
\textbf{Outcome: Conditional reduction.}

The exact catalogue conclusion is shown to fail under the extra
singular-in-$V$/regular-in-HOD configuration, via the cited weak-extender
bridge. The exacting-cardinal literature supplies a conditional source of
that configuration, not an unconditional counterexample from the stated
premise alone. An elementary ultrafilter limit demonstrates a genuine gap
in a naive constructive alternative.

No new relative consistency theorem, unconditional refutation, or
Ultimate-$L$ model is obtained. The remaining task is either a justified
construction under the catalogue premise or a countermodel argument with
its necessary additional consistency assumptions resolved or explicitly
accepted. The notebook's contribution is an audited dependency analysis
and a failed-method test, not priority over the cited work.

% Keep the local bibliography together rather than leave a source fragment.
\needspace{170mm}
% END RESEARCH_ADDITION ATTEMPT_75
\subsection{Sources}
\begin{itemize}\raggedright
\item[\textnormal{[S1]}]\hypertarget{source:30006711:1}{} Bagaria and Ternullo, Steel's Programme: Evidential Framework, the Core and Ultimate-L, Review of Symbolic Logic (2021), Conjecture 5.2.
\url{https://www.cambridge.org/core/journals/review-of-symbolic-logic/article/steels-programme-evidential-framework-the-core-and-ultimatel/22BC5068C2DC48CC899C60E1445A7792}.
{\small\textit{Catalogue formulation source.}}
\item[\textnormal{[S2]}]\hypertarget{source:30006711:2}{} Large cardinals, structural reflection, and the HOD Conjecture.
\url{https://arxiv.org/abs/2411.11568}.
{\small\textit{Catalogue research/status reference; not a proof endorsement.}}
\item[\textnormal{[S3]}]\hypertarget{source:30006711:3}{} The AD+ duality program, the HOD conjecture, and the Ultimate-L conjecture.
\url{https://link.springer.com/content/pdf/10.1007/s11537-026-2541-4.pdf}.
{\small\textit{Catalogue research/status reference; not a proof endorsement.}}
\item[\textnormal{[S4]}]\hypertarget{source:30006711:4}{} Proof of the Omega-Conjecture and Establishment of the V = Ultimate L Axiom: Mathematical Foundations and Physical Realization.
\url{https://www.preprints.org/frontend/manuscript/f77017e8752bef0a5fb615eb22525dd8/download_pub}.
{\small\textit{Catalogue research/status reference; not a proof endorsement.}}
% BEGIN RESEARCH_ADDITION REFERENCES_75
\item[\textnormal{[E1]}]\hypertarget{extra:30006711:1}{} Juan P. Aguilera,
Joan Bagaria and Philipp L\"ucke, \emph{Large cardinals, structural reflection,
and the HOD Conjecture}, arXiv:2411.11568v4, consulted HTML carrying an
internal date of 24 August 2026; Theorems 1.1--1.2, Definition 2.4,
Theorem 2.10 and Conclusion 6.8.
\url{https://arxiv.org/html/2411.11568}.
{\small\textit{Inspected 27 September 2026. The version banner and internal
date are reported separately rather than silently reconciled. The conditional
large-cardinal results are used with their joint hypotheses; their complete
proofs have not been independently certified in this notebook.}}
% END RESEARCH_ADDITION REFERENCES_75
\end{itemize}

\end{document}
